\section{Two approximate coordinate moves}\label{sec:moves}

The quotient coordinates have bounded edge increments, but no symmetry
need interchange their directions. We obtain two directions from
connection probabilities, then normalize them by a linear map. The
interface needed for the corridors is that the seed radius is fixed
before the coordinate increments are made small.

\begin{proposition}[Approximate moves]\label{prop:moves}
Let $G$ be an infinite connected locally finite graph and let a group
$\Lambda$ act on it by automorphisms with finitely many vertex orbits.
Let $\pi:V(G)\to\mathbb Z^2$ be onto, with edge increments
bounded in Euclidean norm by $D_\pi$ and with
$\pi(gx)=\pi(x)+a_g$, where
$g\mapsto a_g$ is an epimorphism from $\Lambda$ to $\mathbb Z^2$.
Suppose $0<\pc<1$ and, at $\pc$, there is almost surely a unique
infinite cluster and no infinite open ray with bounded $\pi$-projection.
For every \(0<e,\eta<1/10\), there is an integer \(m\) such that
for every requested \(\rho_0>0\) one can choose an invertible linear map
\(\Phi:\mathbb R^2\to\mathbb R^2\) and a parameter
\(p'\in[\pc/2,\pc)\) with the following properties.
The coordinates \((t,z)=\Phi\pi\) have edge increments at most
\(\rho_0\) in supremum norm. For $y\in V(G)$ and
$s\in\{-1,1\}$, let $E_s(y)$ be the event that an open path starts
in $B(y,m)$, stays within supremum coordinate distance $100$ of $y$,
and ends within supremum distance $e/4$ of
$(t(y)+1,z(y)+s)$. Then
\[
 \PP_{p'}(E_s(y))\ge1-\eta
 \qquad(y\in V(G),\ s\in\{-1,1\}).
\]
The radius $m$ is chosen before $\rho_0$.
\end{proposition}

We verify the uniqueness hypothesis for the subexponential graph in
our induction before proving the proposition. The adjacent-cluster
bound supplies this fact without any coordinate assumptions.

\begin{lemma}[Uniqueness under subexponential growth]\label{lem:uniqueness}
Let $G$ be infinite, connected, locally finite and quasi-transitive, with
$\log V_+(r)=o(r)$. If $p\in(0,1)$ and Bernoulli bond percolation at $p$
has an infinite cluster with positive probability, then it has exactly
one infinite cluster almost surely.
\end{lemma}

\begin{proof}
Letting $s\to\infty$ in Proposition~\ref{prop:two-cluster} shows that,
for every edge, the probability that its endpoints lie in distinct
infinite clusters is zero. We show that this excludes any two infinite
clusters, even if no edge initially joins them.

Otherwise countability gives fixed vertices $x,y$ that belong to
distinct infinite clusters with positive probability. Choose a fixed
simple path $x=x_0,x_1,\ldots,x_d=y$. On that event let $j$ be the
first index with $x_j\in C_y$; some fixed $j\ge1$ occurs on a
positive-probability subevent. Force the finitely many prefix edges
$x_ix_{i+1}$ for $0\le i<j-1$ open (none if $j=1$). Every endpoint of these edges
was outside $C_y$, so this modification leaves $C_y$ unchanged.
It makes $C_{x_{j-1}}$ infinite through $x$, while that cluster
remains distinct from $C_{x_j}=C_y$. The modified event still has
positive probability: condition on all unmodified bits and use that
the all-open state of the finite prefix has positive conditional
probability. This contradicts the preceding conclusion for the
fixed edge $x_{j-1}x_j$. Thus there is at most one infinite cluster
almost surely. Existence is a tail event, so its assumed positive
probability is one.
\end{proof}

Under the contradiction hypothesis used to prove
Theorem~\ref{thm:structured}, this lemma
and Lemma~\ref{lem:bounded-projection} verify the two percolation
hypotheses of Proposition~\ref{prop:moves}. The equivariant onto
coordinates and their bounded increments were established in
Section~\ref{sec:structure}. The proposition's proof uses only its
stated hypotheses; in particular, absence of bounded-projection rays
is an input supplied by the induction, not by the coordinate map alone.

Choosing geometry by exit probabilities also appears in Martineau and
Tassion's parallelogram construction
\cite[Section~3.2]{MartineauTassion2017}. The proof here obtains the two
directions either from a disk of ellipses or from a projected strip.

\begin{proof}[Proof of Proposition~\ref{prop:moves}]
Choose $o$ with $\pi(o)=0$ and put $\zeta=e/1000$.
Countability gives a vertex with positive infinite-cluster probability.
Harris's inequality transfers this positivity along fixed paths to each
of the finitely many orbit representatives. Hence
\[
 \theta:=\min_x\PP_{\pc}(|C_x|=\infty)>0.
\]
For each fixed nonzero linear functional \(\ell\) on \(\mathbb R^2\),
the unique infinite cluster is unbounded above in \(\ell\circ\pi\),
almost surely. Indeed, the union of all infinite clusters meets
\(\{\ell\circ\pi>r\}\) with probability at least \(\theta\), by
choosing one vertex there. Its unboundedness above is a tail event:
changing finitely many edges changes that union by at most finitely
many vertices. Continuity from above and the zero-one law give
probability one, and uniqueness gives the stated conclusion.

We first construct the two events at \(o\) and at \(\pc\), with error
smaller than \(e/8\) and failure probability smaller than \(\eta/10\).
In each of the two cases below the seed is chosen before the scales
that make the coordinate map small. Only after fixing the map will we
lower the percolation parameter.

\emph{Case 1: no fixed direction supports a strip ray with positive
probability.} For a unit vector $d$, a centered strip means
$\{x:|d^\perp\cdot\pi(x)|\le w\}$ for some finite $w$. Translated
strips are contained in wider centered strips. In this case the
probability of a ray in such a strip is zero for each fixed $d$.
We do not need a statement about the union over all directions.

Cover the unit circle by $M$ arcs of angular radius at most $\zeta$,
with signed unit centers. Choose $m_0$ so that $B(o,m_0)$ hits the
infinite cluster with probability greater than $1-(\eta/20)^M$.
This choice does not involve any later scale. Fix an arbitrarily large
$W$ containing the seed projections well inside the Euclidean ball of
radius $W$. For $U\ge W$, consider the positive matrices
\[
 A(s,\varphi)=R_\varphi
       \begin{pmatrix}W+(U-W)s&0\\0&W\end{pmatrix}R_\varphi^{-1},
 \qquad 0\le s\le1,\quad\varphi\in\mathbb R/\pi\mathbb Z.
\]
The parameter space is a disk: all angles at $s=0$ give $W I$.
Consider open paths from the seed whose last vertex is their first
exit from $\{|A^{-1}\pi|_2<1\}$. Their final normalized radius is
between $1$ and $1+D_\pi/W$. Seed contact with the infinite cluster
guarantees such a path, since an infinite ray cannot have bounded
projection. Split the exit event according to the $M$ endpoint arcs.
Their misses are decreasing events. Harris's inequality implies that
at least one arc is hit with probability greater than $1-\eta/20$.
Call its center a favored center.

For sufficiently large $U$, every favored center at a boundary
parameter $s=1$ has its unoriented line within angle $\pi/16$ of the
minor axis, uniformly in $\varphi$. Otherwise take $U_j\to\infty$
and boundary parameters violating this assertion. A subsequence of
their major axes converges to a fixed direction $d$. Since the arc
radii are small, favored exits have major-axis displacement bounded
below by a positive constant times $U_j$, while the whole exit path
has transverse displacement at most $W+D_\pi$. Loop erasure gives
arbitrarily long simple open paths from the finite seed. For each
fixed length $l$, all such initial segments lie in the finite union
of radius-$l$ balls about the seed. Convergence of the axes is uniform
on the projections of this finite set. Since $D_\pi<W$, it puts
these segments in the centered strip about $d$ of width $2W$, for
all sufficiently large $j$. Such length-$l$ paths thus have probability at least
$1-\eta/20$ for every $l$. These events decrease with $l$, and finite
branching supplies an infinite ray in the strip, a contradiction.

Triangulate the parameter disk so finely that adjacent matrices
satisfy $\|A_i^{-1}A_j-I\|_2\le\zeta$ and boundary angle increments
are less than $\pi/32$. Choose one favored center at each vertex.
Some adjacent center lines have projective separation at least
$\pi/8$. Indeed, otherwise the principal increments of their doubled
line angles have absolute value less than $\pi/4$. Their sum around
each triangle is a multiple of $2\pi$ of absolute value less than
$2\pi$, hence zero; summing triangles gives zero around the boundary.
On the boundary, however, the doubled angles have lifts
$2\varphi+\pi$ with errors of absolute value less than $\pi/8$.
The mesh condition makes successive lift differences principal
increments, whose total is $2\pi$. This is a contradiction.

\begin{figure}[htbp]
\centering
\begin{tikzpicture}[x=1cm,y=1cm,>=Stealth]
 \begin{scope}[shift={(-3.1,0)}]
  \draw[fill=blue!3,draw=blue!50!black] (0,0) circle (1.65);
  \foreach \ang in {0,45,...,315}{
   \draw[black!18] (0,0)--(\ang:1.65);
   \fill[blue!65!black] (\ang:1.65) circle(.035);
   \begin{scope}[shift={(\ang:1.65)},rotate={\ang/2+90}]
    \draw[blue!70!black,line width=1.1pt] (-.19,0)--(.19,0);
   \end{scope}
  }
  \node[fill=white,inner sep=3pt,align=center] at (0,0)
     {$s=0$\\$A=WI$};
  \draw[->,black!60] (15:2) arc[start angle=15,end angle=65,radius=2];
  \node[above] at (0,2.08) {$s=1$};
  \node[below=6pt,align=center,text width=4.9cm] at (0,-1.8)
    {Boundary center lines follow the minor axes.};
 \end{scope}
 \begin{scope}[shift={(2.65,-.05)}]
  \draw[->,black!45] (-.25,0)--(2.35,0) node[right] {$t$};
  \draw[->,black!45] (0,-1.65)--(0,1.7) node[above] {$z$};
  \draw[->,blue!65!black,line width=1pt] (0,0)--(1.65,1.3);
  \draw[->,blue!65!black,line width=1pt] (0,0)--(1.65,-1.3);
  \draw[fill=blue!8,draw=blue!55!black] (1.5,1.15) rectangle (1.8,1.45);
  \draw[fill=blue!8,draw=blue!55!black] (1.5,-1.45) rectangle (1.8,-1.15);
  \fill (0,0) circle(.035);
  \node[right=5pt] at (1.8,1.3) {$(1,1)$};
  \node[right=5pt] at (1.8,-1.3) {$(1,-1)$};
  \node[below=6pt,align=center,text width=4.6cm] at (1,-1.75)
    {A linear map normalizes two separated directions.};
 \end{scope}
\end{tikzpicture}
\caption{The direction argument in Proposition~\ref{prop:moves}.
Left: the ellipse parameters form a disk; on its boundary, favored
unoriented lines have doubled angles with winding number one.
They therefore cannot vary by a small angle across every edge of
a triangulation. Right: two separated favored directions yield
the required moves. The endpoint error boxes are enlarged for clarity.}
\label{fig:moves}
\end{figure}


Take such an adjacent pair with matrices $A_i,A_j$. Let $d_1,d_2$
be their signed favored unit arc centers in their respective normalized
coordinates, so $|\det(d_1,d_2)|\ge\sin(\pi/8)$. Let $\Psi$ send
these original centers to $(1,1)$ and $(1,-1)$; its operator norm is
less than $10$. Set $\Phi=\Psi A_i^{-1}$. A point of the second path
with normalized coordinate $u=A_j^{-1}\pi(x)$ has coordinate
$A_i^{-1}A_j u$ under the first matrix. Its displacement from $u$
is at most $\zeta|u|_2$. Thus arc width, radial overshoot, and this
adjacent-matrix error give endpoint error less than $e/8$ once
$D_\pi/W<\zeta$. The whole paths have coordinate displacement less
than $30$, and edge increments tend to zero as $W\to\infty$.

\emph{Case 2: some fixed direction supports a strip ray with positive
probability.} Existence of a ray in some finite-width centered strip
in that direction is a tail event: a suffix survives every finite
edge modification. Its probability is therefore one. Choose a seed
radius $m_0$ and a width $w$ so that a ray from $B(o,m_0)$ in that
strip exists with probability greater than $1-(\eta/30)^2$.
Enlarge $w$ to contain the seed projection as well. Write the
longitudinal and transverse coordinates in these perpendicular axes.
Such a ray is longitudinally unbounded in absolute value, by the
bounded-projection hypothesis.

Choose an arbitrarily large transverse scale $V$ with
$w/V<\zeta$ and $D_\pi/V<\zeta$.
By uniqueness and the preceding directional tail observation, on
the seed-strip event the seed also connects to absolute transverse
level $V$. Stop at the first such crossing. Its finite path can be
confined to some longitudinal bound $M(V)$ while retaining probability
greater than $1-2(\eta/30)^2$: increase the bound and use continuity
from below. Choose $U$ with $M(V)/U<\zeta$ and $D_\pi/U<\zeta$.
The strip ray also gives a
first crossing of absolute longitudinal level $U$ with probability
greater than $1-(\eta/30)^2$.

Split each kind of exit by its sign. Harris's inequality on the two
misses supplies one favored sign of each kind, each with probability
greater than $1-\eta/10$. After division by $U,V$, each endpoint is
within supremum distance $\zeta$ of its signed coordinate unit vector,
and both paths lie in the square $[-2,2]^2$.
A linear map sending those two vectors to $(1,1),(1,-1)$ gives
endpoint error less than $e/8$ and path displacement less than $30$.
For any requested operator bound, first take $V$ sufficiently large,
then determine the finite bound $M(V)$, and finally take $U$ sufficiently
large. The last linear map has norm at most $2$, so the composite map
has norm at most $2/\min(U,V)$. No bound on the growth of $M(V)$ is
needed. Thus edge increments can be made arbitrarily small after
the seed has been fixed.

Let $c$ be a density radius of the orbit $\Lambda o$ and set
$m=m_0+c$. Given $y$, choose $g$ with $d(go,y)\le c$.
Equivariance translates the two events at $o$ to $go$, and their seeds
lie in $B(y,m)$. Include in the preceding scale choices the requirements
$|\Phi\pi(go)-\Phi\pi(y)|_\infty<e/8$ and that edge increments
are at most the requested $\rho_0$: it suffices to make the operator norm
of $\Phi$ smaller than both $e/(8(c+1)D_\pi)$ and $\rho_0/D_\pi$.
The endpoint errors are then at
most $e/4$, and the path bound is less than $100$. All choices of
$m$ preceded the requested small edge increment.

Only now lower the percolation parameter. Finite paths witness the
two events at $o$. Approximate each by a finite union of path witnesses,
preserving a probability greater than $1-\eta$, with room to spare.
Continuity of these finite events gives the same bounds at a common
$p'\in[\pc/2,\pc)$. Their translates give the required bounds at
every $y$ by the preceding deterministic containment.
\end{proof}
